\section*{Kokkuvõte}\label{kokkuvote} \addcontentsline{toc}{section}{Kokkuvõte}

Bakalaureusetöö eesmärk oli laiendada AKT aines kasutatava õppevahendi \textsc{CMa} virtuaalmasina Java-implementatsiooni. Lisatud sai funktsioonikutsete tugi, tagades täitmiskeskkonna korrektsuse pesastatud ja rekursiivsete väljakutsete korral.

Töös anti ülevaade \textsc{CMa} arhitektuurist, käsustiku semantikast ja funktsioonikutsete rollist programmeerimiskeeltes. Töö tulemusel implementeeriti 13 uut masinakäsku. Selleks laiendati instruktsioonide tüübisüsteemi ning lisati simulaatorile raamipõhise adresseerimise tugi. Lahenduse korrektsus kinnitati 46 JUnit testiga, mis sisaldasid üksuste, integratsiooni- ja regressioonikatsetusi ning testisid muu hulgas rekursiivset faktoriaali ja indekseeritud hüppeprogramme. Kõik testid sooritati edukalt ning ühtegi olemasolevat käsku ei rikutud. Kokku lisati 1251 uut koodirida ning 5 eemaldati.

Käesoleva töö raames dünaamilise mäluhalduse realiseerimist ei käsitletud. Küll aga pakub \textsc{CMa} arhitektuur selleks sobiva lähtekoha ning vastava funktsionaalsuse lisamine oleks loogiline jätk edasises arendustöös, võimaldades modelleerida keerukamaid andmestruktuure ja mälu kasutusmustreid.

Autor ja juhendaja on tehtud tööga rahul: funktsioonikutsete tugi on realiseeritud läbipaistvalt ja spetsifikatsioonipõhiselt ning simulaator on nüüd võimeline modelleerima protseduuriliste programmide juhtvoogu terviklikul kujul.